\documentclass[11pt]{article}
\usepackage[margin=1in]{geometry}
\usepackage{amsmath,amssymb,amsthm}
\usepackage[T1]{fontenc}
\usepackage{lmodern}
\usepackage{microtype}
\usepackage[colorlinks=true,linkcolor=blue,citecolor=blue,urlcolor=blue]{hyperref}
\newtheorem{theorem}{Theorem}
\newtheorem{lemma}[theorem]{Lemma}
\newtheorem{proposition}[theorem]{Proposition}
\newcommand{\C}{\mathbb C}
\newcommand{\R}{\mathbb R}
\newcommand{\Z}{\mathbb Z}
\newcommand{\Q}{\mathbb Q}
\DeclareMathOperator{\Tr}{Tr}
\DeclareMathOperator{\Norm}{N}
\title{Discreteness of sets generated by fixed-parameter extrapolation}
\author{Kaiwen Jin\\\small\href{mailto:yui.ui.siki@gmail.com}{yui.ui.siki@gmail.com}}
\date{October 14, 2026}
\begin{document}
\maketitle
\begin{abstract}
Starting from $0$ and $1$, repeatedly apply the operation
$(a,b)\mapsto(1-\lambda)a+\lambda b$, with a fixed complex parameter
$\lambda$. We prove that the resulting set is discrete exactly when
$\lambda$ is an algebraic integer whose other conjugates, apart from
$\lambda$ and $\overline\lambda$, are real and lie in $(0,1)$.
This gives the converse to the known sufficient condition for separated
extrapolation sets. The proof adjusts integer polynomials within a residue
class while preserving their values at $\lambda$. If there is an additional
conjugate outside $(0,1)$, density of the corresponding field projections
then produces nonzero generated points arbitrarily close to $0$.
\end{abstract}

\section{The discreteness question}

For $\lambda\in\C$, write
\[
a\star_\lambda b=(1-\lambda)a+\lambda b,
\]
and let $Q_\lambda$ be the smallest subset of $\C$ containing $0,1$ and
closed under $\star_\lambda$. Each point of $Q_\lambda$ is obtained by
finitely many operations; no topological closure is taken. The operation
has been used to generate point sets in aperiodic order, where separation
distinguishes discrete patterns from arbitrarily fine ones
\cite{fgh-journal}.

Following Fenner, Green and Homer \cite[Definition 4.1]{fgh-journal},
an algebraic integer $\lambda$ is \emph{strong PV} if every conjugate
other than $\lambda$ and $\overline\lambda$ is real and belongs to
$(0,1)$. Integers and nonreal quadratic algebraic integers satisfy this
condition vacuously. Strong PV parameters give uniformly discrete sets
\cite[Theorem 4.12]{fgh-journal}; discreteness also forces algebraic
integrality \cite[Theorem 8.2]{fgh-extended}. Campbell recently proved
the necessary conjugate restriction for totally real algebraic integers
\cite[Theorem 2]{campbell}. The full complex-parameter converse is
Conjecture 11.1 of \cite{fgh-journal}.

\begin{theorem}\label{thm:main}
For every $\lambda\in\C$, the set $Q_\lambda$ is discrete if and only
if $\lambda$ is strong PV. In that case,
\[
 |x-y|\ge1\qquad(x,y\in Q_\lambda,\ x\ne y).
\]
\end{theorem}

The main issue is the converse. A generated point can be represented by
an integer polynomial that stays between $0$ and $1$ on the real unit
interval. We show that the interval condition can be imposed after fixing
the polynomial's residue modulo the minimal polynomial of $\lambda$.
An extra conjugate outside $(0,1)$ leaves an unbounded field direction;
projecting away that direction gives enough small residues to contradict
discreteness. The argument treats real parameters with nonreal conjugates
and nonreal parameters of arbitrary degree together.

\section{Integer polynomials and prescribed residues}

Let $Q_{[t]}\subseteq\Z[t]$ be generated by $0,1$ under
$(f,g)\mapsto(1-t)f+tg$. Fenner, Green and Homer proved
\cite[Theorem 3.6 and Corollary 3.8]{fgh-journal} that
\begin{align}
Q_{[t]}&=\{0,1\}\cup
 \{P\in\Z[t]:0<P(u)<1\text{ for every }0<u<1\},\label{eq:polynomial}\\
Q_\lambda&=\{P(\lambda):P\in Q_{[t]}\}.\label{eq:evaluation}
\end{align}
These are statements about finite generation. In particular, a polynomial
in the right-hand side of \eqref{eq:polynomial} gives a finite operation
tree for its value at every complex parameter.

Put $q(t)=t(1-t)$. We will need the elementary integer approximation
property behind rounded Bernstein polynomials; see
\cite[Section 1]{draganov} for this classical construction.

\begin{lemma}\label{lem:approximation}
If $h\in C([0,1],\R)$ and $h(0),h(1)\in\Z$, then $h$ can be
uniformly approximated by polynomials $H\in\Z[t]$ with
$H(0)=h(0)$ and $H(1)=h(1)$.
\end{lemma}
\begin{proof}
For $n\ge1$, round each coefficient
$\binom nk h(k/n)$, $1\le k<n$, to an integer $c_k$ with error at
most $1/2$. Set $c_0=h(0)$, $c_n=h(1)$ and
\[
 H_n(t)=\sum_{k=0}^n c_k t^k(1-t)^{n-k}\in\Z[t].
\]
The endpoints are exact. If $B_nh$ is the usual Bernstein polynomial,
\[
 |H_n(t)-B_nh(t)|\le\frac12
   \sum_{k=1}^{n-1}t^k(1-t)^{n-k}.
\]
For $0\le t\le1/4$, the sum is at most
\[
 \frac{t(1-t)^n}{1-2t}\le2t(1-t)^n\le\frac2{n+1}.
\]
The same bound holds on $[3/4,1]$ by symmetry. On $[1/4,3/4]$
the sum is at most $(n-1)(3/4)^n$. Thus the rounding error tends
uniformly to zero.

For completeness, the weights $w_k=\binom nk t^k(1-t)^{n-k}$ satisfy
\[
 \sum_k w_k=1,\qquad
 \sum_k(k/n-t)^2w_k=\frac{t(1-t)}n\le\frac1{4n},
\]
by the binomial identities for the first two moments. If
$\omega_h(\rho)$ denotes the modulus of continuity of $h$, splitting
the sum at $|k/n-t|=\rho$ gives
\[
 |B_nh(t)-h(t)|\le\omega_h(\rho)+
                  \frac{\|h\|_\infty}{2n\rho^2}.
\]
First choose $\rho$ small and then $n$ large. Hence $B_nh\to h$
uniformly, and the same is true of $H_n$.
\end{proof}

\begin{lemma}\label{lem:residue}
Let $m,f\in\Z[t]$ with $m(0)m(1)\ne0$. Suppose that
\[
 0<q(\mu)f(\mu)<1
\]
at every real zero $\mu$ of $m$ in $(0,1)$. There is an $H\in\Z[t]$
such that $P=q(f+mH)$ satisfies $0<P(t)<1$ for all $0<t<1$.
For every zero $\lambda$ of $m$,
\[
 P(\lambda)=q(\lambda)f(\lambda)\in Q_\lambda.
\]
\end{lemma}
\begin{proof}
We first choose a continuous real function $h$ such that
\begin{equation}\label{eq:feasible}
 F(t):=f(t)+m(t)h(t)>0,\qquad q(t)F(t)<1
 \quad(0\le t\le1),
\end{equation}
with integer values of $h$ at the endpoints.
At an interior real zero of $m$, the constant $h=0$ satisfies
\eqref{eq:feasible} on a neighborhood. At an endpoint $e\in\{0,1\}$,
choose $n_e\in\Z$ with $f(e)+m(e)n_e>0$. Since $q(e)=0$, the
constant $h=n_e$ works on a neighborhood of $e$.
At any remaining $t_0\in(0,1)$, the constant
\[
 c_{t_0}=\frac{1/(2q(t_0))-f(t_0)}{m(t_0)}
\]
makes $F(t_0)=1/(2q(t_0))$, so it also works locally.

Choose the interior neighborhoods to avoid the endpoints and each endpoint
neighborhood to avoid the opposite endpoint. A finite subcover of $[0,1]$
has a continuous nonnegative partition of unity. For example, normalize
the distances to the complements of its relative-open members. Take $h$
to be the corresponding convex combination of the local constants.
The inequalities in \eqref{eq:feasible} are strict affine inequalities
in $h(t)$, so they are preserved. At each endpoint only its designated
neighborhood contributes, giving $h(e)=n_e$.

Both continuous functions $F$ and $1-qF$ have positive minima, say
$a$ and $b$. With $M=\max(1,\|m\|_\infty)$, Lemma~\ref{lem:approximation}
gives $H\in\Z[t]$ satisfying
\[
 \|H-h\|_\infty<\frac{\min(a,b)}{2M}.
\]
Then $|m(H-h)|<\min(a,b)/2$. Since $0\le q\le1/4$ on $[0,1]$,
replacing $h$ by $H$ preserves both inequalities. Thus
$0<q(f+mH)<1$ on $(0,1)$. Finally, $m(\lambda)=0$ gives the
stated value, and \eqref{eq:polynomial}--\eqref{eq:evaluation} put it
in the finite generated set $Q_\lambda$.
\end{proof}

\section{Density after omitting a conjugate}

We use a character test for additive subgroups of a real Euclidean space.
Its short proof is included to make the density argument explicit.

\begin{lemma}\label{lem:character}
If an additive subgroup $G\subseteq\R^N$ is not dense, there is a
nonzero real linear functional $L$ with $L(G)\subseteq\Z$.
\end{lemma}
\begin{proof}
Let $A=\overline G$ and choose a vector subspace $V\subseteq A$
of maximal dimension. Orthogonal projection gives
$A=V+D$, where $D=A\cap V^\perp$.
The subgroup $D$ is discrete. Otherwise some nonzero $d_j\in D$
tend to $0$, and a subsequence of $d_j/|d_j|$ tends to a unit
vector $v\in V^\perp$. For every $u\in\R$ choose integers $k_j$
with $k_j|d_j|\to u$. Closedness gives $k_jd_j\to uv\in D$,
so $V+\R v\subseteq A$, contrary to maximality.

Put $W=\operatorname{span}_{\R}D$. If $W\ne V^\perp$, take a
nonzero functional annihilating $V+W$. Otherwise choose a real basis
$d_1,\ldots,d_r$ of $V^\perp$ from $D$ and let
$D_0=\sum_i\Z d_i$. A bounded fundamental parallelepiped for $D_0$
meets the discrete group $D$ in finitely many points, so $D/D_0$ is
finite. Some positive integer $J$ satisfies $JD\subseteq D_0$.
The functional with values $J,0,\ldots,0$ on $d_1,\ldots,d_r$,
extended by zero on $V$, takes integer values on $A$. It is nonzero:
$r>0$ because $A$ is proper. Either case gives the required $L$.
\end{proof}

Let $\lambda$ be an algebraic integer of degree $d$,
$K=\Q(\lambda)$ and $O=\Z[\lambda]$. List the real embeddings
as $\sigma_1,\ldots,\sigma_r$, and choose one representative
$\tau_j$ from each of the $s$ nonreal conjugate pairs. The full
embedding space and map are
\[
 E=\R^r\times\C^s,\qquad r+2s=d,\qquad
 \Sigma(\alpha)=(\sigma_i(\alpha),\tau_j(\alpha)).
\]
A coordinate selection $T$ keeps some real factors and some \emph{whole}
complex factors. Write $\pi_T:E\to E_T$ for the projection.

\begin{proposition}\label{prop:density}
If $T$ is a proper coordinate selection, then $\pi_T\Sigma(O)$
is dense in $E_T$, viewed as a real Euclidean space.
\end{proposition}
\begin{proof}
The empty selection is trivial. Otherwise, if the image is not dense,
Lemma~\ref{lem:character} gives a nonzero real functional $L$ on $E_T$
taking integer values on that image. Let $\widetilde L=L\circ\pi_T$.
Since $1,\lambda,\ldots,\lambda^{d-1}$ is a rational basis of $K$,
nondegeneracy of the trace pairing gives a unique $b\in K$ such that
\[
 \Tr_{K/\Q}(b\lambda^j)=\widetilde L(\Sigma(\lambda^j))
 \quad(0\le j<d).
\]
Indeed, the right-hand sides are integers, and the pairing is nondegenerate
because $b\ne0$ implies $\Tr(b b^{-1})=d\ne0$.

On $E$ this trace is represented by the real functional
\begin{equation}\label{eq:trace}
 T_b((x_i),(z_j))=
 \sum_i\sigma_i(b)x_i+2\operatorname{Re}\sum_j\tau_j(b)z_j.
\end{equation}
The vectors $\Sigma(1),\ldots,\Sigma(\lambda^{d-1})$ form a real
basis of $E$. A real linear relation would give a real polynomial of
degree $<d$ vanishing at all real roots and at one root in every nonreal
pair, hence at their conjugates as well. The $d$ roots of the minimal
polynomial are distinct, so the relation must be zero.

Thus $\widetilde L=T_b$ on $E$. But $\widetilde L$ vanishes on an
omitted coordinate. By \eqref{eq:trace}, the corresponding real or
complex embedding of $b$ is zero; for a complex factor,
$z\mapsto2\operatorname{Re}(cz)$ vanishes identically only if $c=0$.
Field embeddings are injective, so $b=0$ and then $L=0$, a contradiction.
\end{proof}

\section{Proof of the characterization}

\begin{proof}[Proof of Theorem~\ref{thm:main}]
Suppose first that $Q_\lambda$ is discrete. The integrality theorem
\cite[Theorem 8.2]{fgh-extended} makes $\lambda$ an algebraic integer.
If $\lambda=0$ or $1$, it is already strong PV. If $0<\lambda<1$,
the nonzero points $\lambda^n\in Q_\lambda$ tend to $0$, a
contradiction. We may therefore assume $\lambda\notin[0,1]$.

Let $m\in\Z[t]$ be its monic minimal polynomial, and list its real
roots in $(0,1)$ as $\mu_1,\ldots,\mu_k$. Irreducibility gives
$m(0)m(1)\ne0$. In the full embedding space retain the physical
coordinate corresponding to $\lambda$, together with all these $k$
real coordinates. Choose the representative of the physical complex pair
to be $\lambda$ itself when $\lambda$ is nonreal. The resulting space is
\[
 E_T=F\times\R^k,\qquad
 F=\begin{cases}\R,&\lambda\in\R,\\\C,&\lambda\notin\R.\end{cases}
\]

If $\lambda$ is not strong PV, a conjugate other than
$\lambda,\overline\lambda$ lies outside the real interval $(0,1)$.
It belongs to an omitted real factor or an omitted complex pair. Hence
$T$ is proper, and Proposition~\ref{prop:density} applies.
For any $\varepsilon>0$, density gives $\alpha=f(\lambda)\in O$,
$f\in\Z[t]$, with its selected coordinates in the nonempty open set
\begin{equation}\label{eq:window}
 \left\{z\in F:\frac{\varepsilon}{2|q(\lambda)|}<|z|
                    <\frac{\varepsilon}{|q(\lambda)|}\right\}
 \times\prod_{i=1}^k
 \left(\frac1{4q(\mu_i)},\frac3{4q(\mu_i)}\right).
\end{equation}
Here $q(\lambda)\ne0$; the empty product is allowed.
Thus $\alpha\ne0$, $|q(\lambda)\alpha|<\varepsilon$, and
$0<q(\mu_i)f(\mu_i)<1$ for every $i$.
Lemma~\ref{lem:residue} puts the nonzero number
$q(\lambda)\alpha$ in $Q_\lambda$. Since $\varepsilon$ was arbitrary,
$0$ is an accumulation point, contradicting discreteness. This proves
the necessary conjugate restriction.

Conversely let $\lambda$ be strong PV and take distinct
$x,y\in Q_\lambda$. Tree induction gives $a=x-y\in\Z[\lambda]$.
At every other real conjugate $\mu\in(0,1)$ the evaluated trees
stay in $[0,1]$, so $|\sigma_\mu(a)|\le1$. The norm of the nonzero
algebraic integer $a$ is a nonzero integer. Its physical factor is
$a$ in the real case and $a\overline a$ in the nonreal case. Therefore
\[
 1\le|\Norm_{K/\Q}(a)|\le
 \begin{cases}|a|,&\lambda\in\R,\\|a|^2,&\lambda\notin\R.\end{cases}
\]
This gives $|x-y|\ge1$. For $\lambda=0,1$ the set is exactly
$\{0,1\}$, so the same separation holds.
\end{proof}

\paragraph{AI use.}
The proof and exposition were developed with assistance from OpenAI Codex,
including source retrieval and checks of the argument. Those preparation
checks do not constitute independent human review. Independent mathematical
review of the submitted argument is requested.

\begin{thebibliography}{9}
\bibitem{fgh-journal}
S.~Fenner, F.~Green and S.~Homer,
\emph{Fixed-Parameter Extrapolation and Aperiodic Order},
Discrete Comput. Geom. \textbf{76} (2026), 1--74.
\href{https://doi.org/10.1007/s00454-025-00816-4}{doi:10.1007/s00454-025-00816-4}.

\bibitem{fgh-extended}
S.~Fenner, F.~Green and S.~Homer,
\emph{Fixed-Parameter Extrapolation and Aperiodic Order}, extended manuscript,
\href{https://arxiv.org/abs/1212.2889v6}{arXiv:1212.2889v6} (2020).

\bibitem{campbell}
J.~M.~Campbell, \emph{Pinch's conjecture on $a$-convexity},
\href{https://arxiv.org/abs/2609.30771v1}{arXiv:2609.30771v1} (2026).

\bibitem{draganov}
B.~R.~Draganov,
\emph{Simultaneous approximation by Bernstein polynomials with integer coefficients},
\href{https://arxiv.org/abs/1804.08248v2}{arXiv:1804.08248v2} (2018).
\end{thebibliography}
\end{document}
